\chapter{Securities, Trading, and Market Conventions}

\begin{objectives}
Understand the contractual cash flows of major instruments, translate market quotes
into economic exposures, and recognize how settlement, collateral, and conventions
enter valuation.
\end{objectives}

\section{The balance-sheet view}

Every security is simultaneously an asset to one party and a liability or residual
claim of another. A useful first analysis lists:

\begin{itemize}
\item contractual cash flows and their priority;
\item states in which cash flows change;
\item collateral, margin, and close-out rules;
\item currency, settlement date, day count, and calendar;
\item liquidity, financing, and counterparty exposure.
\end{itemize}

Prices are not cash flows. A model maps uncertain future cash flows and discounting
to a price under explicit assumptions.

\section{Equity}

Common equity is a residual claim after contractual obligations. Corporate actions
include dividends, splits, repurchases, rights issues, mergers, and spin-offs.
Raw closing prices across these events do not produce total returns. If $P_t$ is
ex-dividend price and $D_t$ cash distribution,
\[
R_t=\frac{P_t+D_t}{P_{t-1}}-1.
\]
Adjusted-price vendors encode corporate actions with methodology that must be
documented.

Market capitalization weights are proportional to float-adjusted value. Equal
weighting, fundamental weighting, and price weighting create distinct rebalancing
and factor exposures. Index membership is governed by rules and changes over time;
using today's constituents historically creates survivorship bias.

\section{Money-market instruments}

Treasury bills, commercial paper, certificates of deposit, repo, and overnight
indexed exposures use different quotation bases. A discount yield can be based on
face value and a 360-day year; an investment yield uses purchase price. Comparing
numbers without converting conventions creates false arbitrage.

In a repo, one party sells collateral and agrees to repurchase it. The repo rate is
the financing rate; haircut protects the cash lender from collateral-value changes.
Special collateral can finance below the general collateral rate because possession
of a particular security has value.

\section{Bonds}

A fixed-rate bond pays coupons and principal subject to credit and contractual terms.
Clean price excludes accrued interest; dirty price is the settlement amount:
\[
P_{\mathrm{dirty}}=P_{\mathrm{clean}}+\text{accrued interest}.
\]
Yield to maturity is the single internal rate equating discounted promised cash
flows to price. It is not an expected return unless coupons reinvest at that yield,
the bond is held to maturity, and promised cash flows occur.

Floating-rate notes reset coupons to a reference plus spread. Inflation-linked bonds
index principal or coupons with publication lags and floors. Callable bonds embed an
option held by the issuer; putable bonds embed one held by the investor.

\section{Forwards and futures}

A long forward with delivery price $K$ has payoff $S_T-K$. With no income and
deterministic rate, no arbitrage gives $K=S_0\e^{rT}$. For an asset with known
continuous yield $q$,
\[
F_{0,T}=S_0\e^{(r-q)T}.
\]
Storage cost and convenience yield modify commodity carry.

Futures are marked to market daily. When rates are stochastic and correlated with
the underlying, futures and forward prices can differ because gains and losses are
realized at different times.

\section{Swaps}

An interest-rate swap exchanges fixed and floating cash flows on a notional. At
inception, the par fixed rate makes value zero:
\[
K_{\mathrm{swap}}=
\frac{P(0,T_0)-P(0,T_n)}
{\sum_{i=1}^n\delta_iP(0,T_i)}
\]
under a single-curve abstraction. Modern collateralized valuation often projects
floating cash flows on a tenor curve and discounts on an overnight curve.

Currency swaps exchange principal and interest in different currencies. Total-return
swaps transfer the price and income performance of an asset against financing.
Contract details determine funding, collateral, and counterparty risk.

\section{Options}

A European call pays $(S_T-K)^+$ and a put $(K-S_T)^+$. American exercise is
permitted before maturity; path-dependent options use the history of the underlying.
At fixed maturity, call value is decreasing and convex in strike under no static
arbitrage.

\begin{proposition}[Put--call parity]
For European options with the same $K,T$ on a non-dividend-paying asset,
\[
C_0-P_0=S_0-KP(0,T).
\]
\end{proposition}
\begin{proof}
Both portfolios---long call plus the present value of $K$, and long put plus the
asset---pay $\max(S_T,K)$ at $T$. No arbitrage equates their values.
\end{proof}

With known dividends, replace spot by prepaid forward value. Violations after bid,
ask, funding, settlement, shorting, and exercise details are needed before claiming
an executable arbitrage.

\section{Problems}

\begin{exercise}[Core: carry]
Derive equity-index forward price with discrete known dividends and with a continuous
dividend yield. State the assumptions.
\end{exercise}
\begin{exercise}[Core: swap rate]
Replicate a floating-rate note to derive the par swap rate formula.
\end{exercise}
\begin{exercise}[Advanced: futures convexity]
Use a two-period tree with stochastic rates to construct a case in which forward and
futures prices differ.
\end{exercise}


\chapter{Market Microstructure and Execution}

\begin{objectives}
Model the limit order book, spreads, information, impact, and execution; translate a
statistical signal into a feasible trading rule with capacity and timestamp
discipline.
\end{objectives}

\section{Orders and the limit order book}

A market order demands immediate execution at available quotes. A limit order sets
a worst acceptable price and supplies liquidity while exposed to non-execution and
adverse selection. The best bid $b_t$ and ask $a_t$ define quoted midpoint
$m_t=(a_t+b_t)/2$ and spread $s_t=a_t-b_t$.

The transaction price can be written
\[
p_t=m_t+q_t\frac{s_t}{2}+\eta_t,
\]
where $q_t=+1$ for buyer initiation and $-1$ for seller initiation, with $\eta_t$
capturing price improvement and classification error. Alternation between bid and
ask creates negative first-order autocorrelation even when the efficient price is a
martingale.

\section{Why spreads exist}

Spread compensates order-processing cost, inventory risk, and adverse selection.
An informed trader preferentially buys before price rises and sells before it falls,
so passive quotes lose conditionally to informed flow. Dealers widen or adjust
quotes as the probability of information increases.

Effective spread for a buy at price $p_t$ relative to pre-trade midpoint $m_t$ is
$2(p_t-m_t)$. Realized spread compares execution with a later midpoint and attempts
to subtract information-driven price movement. Price impact is the complementary
midpoint response.

\section{Order flow and price discovery}

Signed order flow can predict short-horizon price changes because trades reveal
information and mechanically consume depth. Long memory in trade signs need not
imply easily exploitable return predictability: impact can be transient or
state-dependent and execution costs absorb it.

Hasbrouck-style vector autoregressions, Kyle's linear model, and structural
microstructure models quantify different notions of information and impact. They
rely on timestamp alignment and trade classification that vary across venues.

\begin{example}[Kyle impact]
In a one-period Gaussian model, an informed trader submits $x=\beta(v-p_0)$ and noise
traders submit $u$. A competitive market maker sets
$p=\E[v\mid x+u]=p_0+\lambda(x+u)$. Linear equilibrium chooses $\beta$ and
$\lambda$ so the informed trader optimizes profit and the market maker breaks even.
Impact $\lambda$ rises with fundamental uncertainty and falls with noise-trading
scale.
\end{example}

\section{Transaction-cost models}

A simple cost for trade $\Delta w$ is
\[
C(\Delta w)=
\underbrace{\sum_i c_i\abs{\Delta w_i}}_{\text{spread and fees}}
+\underbrace{\sum_i k_i\abs{\Delta w_i}^{3/2}}_{\text{nonlinear impact}}.
\]
The square-root impact rule is an empirical approximation, not universal. Participation
rate, volatility, volume, urgency, venue, and order type matter.

Portfolio turnover can be defined as
$\frac12\sum_i\abs{w_{i,t}^{\mathrm{target}}-w_{i,t}^{\mathrm{pretrade}}}$.
State which convention is used. Costs must apply to dollar trades after accounting
for drift in pre-trade weights.

\section{Optimal execution}

In the Almgren--Chriss framework, a trader liquidates inventory $x(t)$ by balancing
expected temporary impact against variance from holding risk. A continuous
objective is
\[
\min_{x(0)=X,\ x(T)=0}
\int_0^T\left(\eta\dot x(t)^2+
\lambda\sigma^2x(t)^2\right)\dd t.
\]
The Euler--Lagrange equation is
$\eta\ddot x=\lambda\sigma^2x$, producing a hyperbolic-sine schedule. Larger risk
aversion or volatility accelerates trading; greater temporary impact slows it.

Real execution adds stochastic liquidity, price limits, queues, hidden orders,
fill uncertainty, multiple venues, and feedback between child orders and signals.

\section{Timestamps and backtests}

A research row should distinguish:

\begin{itemize}
\item event time when the economic event occurs;
\item publication time when a vendor makes data available;
\item ingestion time when the strategy receives it;
\item decision time when the signal is frozen;
\item execution window and settlement.
\end{itemize}

Trading at the close with a signal computed from that same official close is usually
impossible unless an executable auction mechanism and submission cutoff are modeled.
Daily bar data hide the path and queue within the bar.

\section{Capacity and crowding}

Capacity is the capital at which expected marginal alpha equals marginal costs and
constraints. Scaling leverage leaves returns unchanged only in a frictionless
linear model. Impact, borrow availability, margin, concentration, and correlated
liquidation make strategy P\&L nonlinear in capital.

Crowding matters when many strategies share factors, signals, or risk controls.
During deleveraging, normally diverse positions can become correlated through common
flows. Stress tests should include spread widening, volume collapse, borrow recall,
and delayed execution.

\section{Market integrity}

Research must distinguish legitimate liquidity provision from manipulative intent.
Spoofing, layering, wash trading, misuse of confidential order information, and
trading on material nonpublic information raise legal and ethical issues beyond
model performance. Market-data licenses can restrict storage, redistribution, and
publication.

\section{Problems}

\begin{exercise}[Core: bid--ask bounce]
Assume a constant efficient price and independent trade signs. Compute the
first-order autocovariance of transaction-price changes.
\end{exercise}
\begin{exercise}[Core: optimal execution]
Solve the Euler--Lagrange equation for the stated boundary conditions and examine
the risk-neutral limit $\lambda\downarrow0$.
\end{exercise}
\begin{exercise}[Research]
Design a transaction-cost calibration using order or bar data. Specify target,
features, sample split, weighting, and how you would handle endogenous order size.
\end{exercise}


\chapter{Arbitrage, Collateral, and Financial Engineering Practice}

\begin{objectives}
Use replication and dominance arguments, understand collateralized valuation,
recognize counterparty and funding adjustments, and establish a professional model
governance process.
\end{objectives}

\section{Law of one price and arbitrage}

If two portfolios deliver identical cash flows in every state and have identical
institutional treatment, they must have the same price in an ideal market. An
arbitrage is a self-financing strategy with zero or negative initial cost,
nonnegative terminal value almost surely, and positive value with positive
probability.

The phrase identical institutional treatment matters. Assets can differ in
liquidity, financing eligibility, taxes, capital charges, haircuts, legal netting,
or delivery optionality. A price basis can compensate these differences without
being a free lunch.

\section{State prices and complete markets}

In a finite-state one-period economy, let payoff matrix $X\in\R^{S\times N}$ and
price vector $p$. A positive state-price vector $q>0$ prices assets when
$p=X^\transpose q$. Absence of arbitrage is equivalent to existence of such a
strictly positive linear pricing functional under standard finite-dimensional
conditions.

The market is complete if $\rank(X)=S$: every state-contingent payoff is replicable
and state prices are unique. In incomplete markets, arbitrage-free prices of an
unspanned claim form an interval or depend on an added criterion.

\begin{theorem}[Finite-state fundamental theorem]
In a finite one-period market, no arbitrage holds if and only if there exists
$q\in\R^S$ with $q_s>0$ and $p=X^\transpose q$.
\end{theorem}
\begin{proof}[Proof sketch]
If positive state prices exist, any nonnegative nonzero payoff has positive cost, so
no arbitrage exists. Conversely, separation of the attainable-payoff cone from a
strictly positive payoff produces a positive linear functional agreeing with market
prices.
\end{proof}

\section{Collateral and discounting}

A collateral agreement specifies eligible assets, currency, threshold, minimum
transfer amount, frequency, haircut, and remuneration. If variation margin is
continuous, cash, and remunerated at overnight rate $c_t$, a stylized collateralized
payoff is discounted at $c$:
\[
V_t=\E_t^\Q\left[
\e^{-\int_t^Tc_s\dd s}H_T
\right].
\]
Changing collateral currency can change value through funding and cross-currency
basis. Optional collateral currency is itself an option.

\section{Counterparty credit and valuation adjustments}

Positive future exposure creates loss if the counterparty defaults. Expected
positive exposure at $t$ is $\E[V_t^+]$; potential future exposure is a quantile.
A stylized unilateral CVA is
\[
\mathrm{CVA}\approx
(1-R)\int_0^T
\E^\Q[D(0,t)V_t^+\mid\tau=t]\,\dd\Prob^\Q(\tau\leq t).
\]
Wrong-way risk occurs when exposure rises as counterparty credit quality deteriorates.

Modern valuation adjustment language also includes DVA for own default, FVA for
funding, MVA for initial margin, and KVA for capital. Summing independently computed
adjustments can double-count effects; the institution needs a consistent replication
and accounting framework.

\section{Hedging as a control problem}

A model price is only one output. A desk holds a hedging policy under discrete
trading, transaction costs, jumps, liquidity limits, and parameter uncertainty.
Hedging P\&L decomposes into model theta, realized gamma, volatility repricing,
cross-Greeks, carry, funding, and residual.

For a delta-hedged option over a short interval under continuous diffusion,
\[
\Delta\Pi\approx
\frac12\Gamma S^2
(\sigma_{\mathrm{real}}^2-\sigma_{\mathrm{imp}}^2)\Delta t
\]
under a simplified convention. The formula motivates variance trading but omits
jumps, surface dynamics, costs, and discrete hedge error.

\section{Model lifecycle}

A controlled model process contains:

\begin{enumerate}
\item purpose, owner, users, and prohibited uses;
\item mathematical specification and assumptions;
\item data lineage and implementation tests;
\item calibration and benchmark comparison;
\item sensitivity, stress, and limitation analysis;
\item independent validation and approval;
\item monitoring thresholds and change control;
\item retirement or remediation criteria.
\end{enumerate}

Verification asks whether code correctly implements the specification. Validation
asks whether the specification is adequate for its use. A correct implementation of
an unsuitable model is still model risk.

\section{Reproducibility and auditability}

Every reported number should be traceable to immutable raw data or an identified
vendor query, code version, configuration, dependency environment, and generation
timestamp. Random seeds support debugging but do not create statistical validity.
Figures should expose units, sampling frequency, and whether returns include
distributions.

\begin{researchnote}
A strong quant portfolio is not a gallery of final Sharpe ratios. It is a chain of
questions, mathematical arguments, failed specifications, robustness checks, and
reproducible decisions.
\end{researchnote}

\section{Problems}

\begin{exercise}[Core: state prices]
For two states and a bond plus stock, solve state prices and price a call. Determine
when the market is complete.
\end{exercise}
\begin{exercise}[Core: CVA]
Explain the difference between exposure under the physical and pricing measures.
Which belongs in pricing and which in risk-management scenario design?
\end{exercise}
\begin{exercise}[Advanced: governance]
Write a validation plan for an illiquid-option model. Include benchmarks, data
limitations, hedging tests, uncertainty, and usage restrictions.
\end{exercise}
